
\documentclass{article}
\usepackage{colortbl}
\usepackage{makecell}
\usepackage{multirow}
\usepackage{supertabular}

\begin{document}

\newcounter{utterance}

\centering \large Interaction Transcript for game `clean\_up', experiment `2\_hard\_3obj\_mt', episode 1 with glm{-}5.2.
\vspace{24pt}

{ \footnotesize  \setcounter{utterance}{1}
\setlength{\tabcolsep}{0pt}
\begin{supertabular}{c@{$\;$}|p{.15\linewidth}@{}p{.15\linewidth}p{.15\linewidth}p{.15\linewidth}p{.15\linewidth}p{.15\linewidth}}
    \# & $\;$A & \multicolumn{4}{c}{Game Master} & $\;\:$B\\
    \hline

    \theutterance \stepcounter{utterance}  
    & & \multicolumn{4}{p{0.6\linewidth}}{
        \cellcolor[rgb]{0.9,0.9,0.9}{
            \makecell[{{p{\linewidth}}}]{
                \texttt{\tiny{[P1$\langle$GM]}}
                \texttt{Jiena l{-}kap tal{-}logħba tiegħek, u inti qed tilgħab logħba kollaborattiva bil{-}grilja li ġejja bħala l{-}bord tal{-}logħba:} \\
\\ 
\texttt{```} \\
\texttt{ 1234567} \\
\texttt{╔═══╤═╤═╗} \\
\texttt{║◌◌◌│◌│◌║ 1} \\
\texttt{║◌P◌└─┤◌║ 2} \\
\texttt{║◌◌◌C◌├─╢ 3} \\
\texttt{╟─┐◌┌─┤◌║ 4} \\
\texttt{║◌│◌│◌│◌║ 5} \\
\texttt{╟─┼─┼─┼─╢ 6} \\
\texttt{║◌│◌│◌│L║ 7} \\
\texttt{╚═╧═╧═╧═╝} \\
\\ 
\texttt{```} \\
\\ 
\texttt{* It{-}tarf ta' fuq juri l{-}koordinati x li jiżdiedu lejn il{-}lemin, u t{-}tarf tal{-}lemin il{-}koordinati y li jiżdiedu 'l isfel.} \\
\texttt{* L{-}oġġetti li ġejjin huma mqiegħda b'mod każwali fuq il{-}grilja tiegħek: 'C', 'L', 'P'.} \\
\\ 
\\ 
\texttt{Il{-}plejer l{-}ieħor jara varjazzjoni tal{-}bord tal{-}logħba, fejn l{-}oġġetti huma mqiegħda f'postijiet każwali differenti. Int ma tistax tara l{-}bord tal{-}plejer l{-}ieħor, u lanqas hu/hi ma jista'/tista' jara tiegħek.} \\
\\ 
\\ 
\texttt{**Għan:**} \\
\\ 
\texttt{Iż{-}żewġ plejers iridu jċaqalqu l{-}oġġetti fuq l{-}isfond rispettiv tagħhom sabiex oġġetti identiċi jispiċċaw fl{-}istess koordinati. Trid tikkomunika mal{-}plejer l{-}ieħor biex taqblu fuq konfigurazzjoni komuni tal{-}għan.} \\
\\ 
\\ 
\texttt{**Regoli:**} \\
\\ 
\texttt{* F'kull dawra, tista' tibgħat eżattament wieħed miż{-}żewġ kmandi li ġejjin:} \\
\texttt{ 1. `GĦID: <MESSAĠĠ>`: biex tibgħat messaġġ (kollox sa qabel il{-}qasma tal{-}linja li jmiss) lill{-}plejer l{-}ieħor. Jiena ngħaddih lis{-}sieħeb/sieħba tiegħek.} \\
\texttt{2. `ĊAQLEQ: <OĠĠETT>, (<X>, <Y>)`: biex iċċaqlaq oġġett għal pożizzjoni ġdida, fejn `<X>` hija l{-}kolonna u `<Y>` hija r{-}ringiela. Ngħidlek jekk iċ{-}ċaqliq tiegħek kienx ta' suċċess jew le.} \\
\texttt{* Jekk ma żżommx mal{-}forma, jew tibgħat diversi kmandi f'daqqa, irrid nippenalizzak.} \\
\texttt{* Jekk iż{-}żewġ plejers jakkumulaw aktar minn 8 penalitajiet, it{-}tnejn titlef il{-}logħba.} \\
\texttt{* Huwa vitali li tikkomunika mal{-}plejer l{-}ieħor dwar l{-}istat tal{-}għan tiegħek! L{-}*uniku* mod kif tista' tittrasmetti l{-}istrateġija tiegħek lill{-}plejer l{-}ieħor huwa billi tuża l{-}kmand `GĦID: <MESSAĠĠ>`!} \\
\\ 
\\ 
\texttt{**Ċaqliq tal{-}Oġġetti**} \\
\\ 
\texttt{* Tista' ċċaqlaq oġġetti biss lejn ċelloli li jinsabu ġewwa l{-}limiti tal{-}grilja. Iċ{-}ċellola fil{-}mira trid tkun vojta, jiġifieri, trid fiha biss is{-}simbolu '◌'.} \\
\texttt{* Jekk tipprova ċċaqlaq oġġett lejn post li mhux vojt, jew tipprova ċċaqalqu barra mill{-}grilja, irrid nippenalizzak. Ikollok prova oħra.} \\
\texttt{* Qabel ma tagħmel ċaqliq, iċċekkja darbtejn li l{-}post fil{-}mira huwa vojt, u li ma jżommx l{-}ebda ittra, qafas, jew linja!} \\
\\ 
\\ 
\texttt{**Tmiem il{-}Logħba**} \\
\\ 
\texttt{Jekk taħseb li wasalt għall{-}għan li tallinjaw l{-}oġġetti kollha, tista' titlob lill{-}plejer l{-}ieħor biex itemm il{-}logħba billi tibgħat `GĦID: lestejt?`. Jekk il{-}plejer l{-}ieħor jitlobek biex itemm il{-}logħba, u inti twieġeb b'`GĦID: lestejt!`, il{-}logħba tintemm.} \\
\\ 
\texttt{Iż{-}żewġ plejers jirbħu jekk il{-}logħba tispiċċa fi żmien 12 rawnd, fejn rawnd wieħed huwa definit bħala żewġ plejers li kull wieħed jibgħat kmand validu.} \\
\\ 
\texttt{**Punteġġ:**} \\
\\ 
\texttt{Iktar ma l{-}oġġetti identiċi jkunu qrib xulxin fiż{-}żewġ bordijiet tal{-}logħba, iktar taqla' punti. Il{-}penalitajiet inaqqsu l{-}punti tiegħek. Tista' tgħaddi r{-}rekord?} \\
\\ 
\\ 
\texttt{Jekk jogħġbok ibgħat messaġġ lill{-}plejer l{-}ieħor biex tibda l{-}logħba!} \\
            }
        }
    }
    & \\ \\

    \theutterance \stepcounter{utterance}  
    & \multicolumn{4}{p{0.6\linewidth}}{
        \cellcolor[rgb]{0.9,0.9,0.9}{
            \makecell[{{p{\linewidth}}}]{
                \texttt{\tiny{[P1$\rangle$GM]}}
                \texttt{GĦID: Hello! Biex nallinejaw l{-}oġġetti tagħna (C, L, P) bl{-}aħjar mod possibbli, nipproponi li nikkonsensus fuq pożizzjonijiet fissi bħala għan komuni. Nipproponi npoġġu l{-}oġġetti f'dawn il{-}koordinati eżatti: C f' (1, 4), L f' (7, 7), u P f' (3, 5). Jekk taqbel, jiena bdejt niċċaqlaq l{-}oġġetti tiegħi lejn dawn il{-}posizzjonijiet.} \\
            }
        }
    }
    & & \\ \\

    \theutterance \stepcounter{utterance}  
    & & & \multicolumn{2}{p{0.3\linewidth}}{
        \cellcolor[rgb]{0.9,0.9,0.9}{
            \makecell[{{p{\linewidth}}}]{
                \texttt{\tiny{[GM$|$GM]}}
                \texttt{GĦID: Hello! Biex nallinejaw l{-}oġġetti tagħna (C, L, P) bl{-}aħjar mod possibbli, nipproponi li nikkonsensus fuq pożizzjonijiet fissi bħala għan komuni. Nipproponi npoġġu l{-}oġġetti f'dawn il{-}koordinati eżatti: C f' (1, 4), L f' (7, 7), u P f' (3, 5). Jekk taqbel, jiena bdejt niċċaqlaq l{-}oġġetti tiegħi lejn dawn il{-}posizzjonijiet.} \\
            }
        }
    }
    & & \\ \\

    \theutterance \stepcounter{utterance}  
    & & \multicolumn{4}{p{0.6\linewidth}}{
        \cellcolor[rgb]{0.9,0.9,0.9}{
            \makecell[{{p{\linewidth}}}]{
                \texttt{\tiny{[GM$\rangle$P2]}}
                \texttt{Jiena l{-}kap tal{-}logħba tiegħek, u inti qed tilgħab logħba kollaborattiva bil{-}grilja li ġejja bħala l{-}bord tal{-}logħba:} \\
\\ 
\texttt{```} \\
\texttt{ 1234567} \\
\texttt{╔═══╤═╤═╗} \\
\texttt{║◌◌◌│◌│◌║ 1} \\
\texttt{║◌◌L└─┤◌║ 2} \\
\texttt{║◌◌◌◌◌├─╢ 3} \\
\texttt{╟─┐◌┌─┤◌║ 4} \\
\texttt{║C│◌│◌│◌║ 5} \\
\texttt{╟─┼─┼─┼─╢ 6} \\
\texttt{║P│◌│◌│◌║ 7} \\
\texttt{╚═╧═╧═╧═╝} \\
\\ 
\texttt{```} \\
\\ 
\texttt{* It{-}tarf ta' fuq juri l{-}koordinati x li jiżdiedu lejn il{-}lemin, u t{-}tarf tal{-}lemin il{-}koordinati y li jiżdiedu 'l isfel.} \\
\texttt{* L{-}oġġetti li ġejjin huma mqiegħda b'mod każwali fuq il{-}grilja tiegħek: 'C', 'L', 'P'.} \\
\\ 
\\ 
\texttt{Il{-}plejer l{-}ieħor jara varjazzjoni tal{-}bord tal{-}logħba, fejn l{-}oġġetti huma mqiegħda f'postijiet każwali differenti. Int ma tistax tara l{-}bord tal{-}plejer l{-}ieħor, u lanqas hu/hi ma jista'/tista' jara tiegħek.} \\
\\ 
\\ 
\texttt{**Għan:**} \\
\\ 
\texttt{Iż{-}żewġ plejers iridu jċaqalqu l{-}oġġetti fuq l{-}isfond rispettiv tagħhom sabiex oġġetti identiċi jispiċċaw fl{-}istess koordinati. Trid tikkomunika mal{-}plejer l{-}ieħor biex taqblu fuq konfigurazzjoni komuni tal{-}għan.} \\
\\ 
\\ 
\texttt{**Regoli:**} \\
\\ 
\texttt{* F'kull dawra, tista' tibgħat eżattament wieħed miż{-}żewġ kmandi li ġejjin:} \\
\texttt{ 1. `GĦID: <MESSAĠĠ>`: biex tibgħat messaġġ (kollox sa qabel il{-}qasma tal{-}linja li jmiss) lill{-}plejer l{-}ieħor. Jiena ngħaddih lis{-}sieħeb/sieħba tiegħek.} \\
\texttt{2. `ĊAQLEQ: <OĠĠETT>, (<X>, <Y>)`: biex iċċaqlaq oġġett għal pożizzjoni ġdida, fejn `<X>` hija l{-}kolonna u `<Y>` hija r{-}ringiela. Ngħidlek jekk iċ{-}ċaqliq tiegħek kienx ta' suċċess jew le.} \\
\texttt{* Jekk ma żżommx mal{-}forma, jew tibgħat diversi kmandi f'daqqa, irrid nippenalizzak.} \\
\texttt{* Jekk iż{-}żewġ plejers jakkumulaw aktar minn 8 penalitajiet, it{-}tnejn titlef il{-}logħba.} \\
\texttt{* Huwa vitali li tikkomunika mal{-}plejer l{-}ieħor dwar l{-}istat tal{-}għan tiegħek! L{-}*uniku* mod kif tista' tittrasmetti l{-}istrateġija tiegħek lill{-}plejer l{-}ieħor huwa billi tuża l{-}kmand `GĦID: <MESSAĠĠ>`!} \\
\\ 
\\ 
\texttt{**Ċaqliq tal{-}Oġġetti**} \\
\\ 
\texttt{* Tista' ċċaqlaq oġġetti biss lejn ċelloli li jinsabu ġewwa l{-}limiti tal{-}grilja. Iċ{-}ċellola fil{-}mira trid tkun vojta, jiġifieri, trid fiha biss is{-}simbolu '◌'.} \\
\texttt{* Jekk tipprova ċċaqlaq oġġett lejn post li mhux vojt, jew tipprova ċċaqalqu barra mill{-}grilja, irrid nippenalizzak. Ikollok prova oħra.} \\
\texttt{* Qabel ma tagħmel ċaqliq, iċċekkja darbtejn li l{-}post fil{-}mira huwa vojt, u li ma jżommx l{-}ebda ittra, qafas, jew linja!} \\
\\ 
\\ 
\texttt{**Tmiem il{-}Logħba**} \\
\\ 
\texttt{Jekk taħseb li wasalt għall{-}għan li tallinjaw l{-}oġġetti kollha, tista' titlob lill{-}plejer l{-}ieħor biex itemm il{-}logħba billi tibgħat `GĦID: lestejt?`. Jekk il{-}plejer l{-}ieħor jitlobek biex itemm il{-}logħba, u inti twieġeb b'`GĦID: lestejt!`, il{-}logħba tintemm.} \\
\\ 
\texttt{Iż{-}żewġ plejers jirbħu jekk il{-}logħba tispiċċa fi żmien 12 rawnd, fejn rawnd wieħed huwa definit bħala żewġ plejers li kull wieħed jibgħat kmand validu.} \\
\\ 
\texttt{**Punteġġ:**} \\
\\ 
\texttt{Iktar ma l{-}oġġetti identiċi jkunu qrib xulxin fiż{-}żewġ bordijiet tal{-}logħba, iktar taqla' punti. Il{-}penalitajiet inaqqsu l{-}punti tiegħek. Tista' tgħaddi r{-}rekord?} \\
\\ 
\texttt{Il{-}plejer l{-}ieħor beda l{-}logħba billi bagħat dan il{-}messaġġ:} \\
\\ 
\texttt{"Hello! Biex nallinejaw l{-}oġġetti tagħna (C, L, P) bl{-}aħjar mod possibbli, nipproponi li nikkonsensus fuq pożizzjonijiet fissi bħala għan komuni. Nipproponi npoġġu l{-}oġġetti f'dawn il{-}koordinati eżatti: C f' (1, 4), L f' (7, 7), u P f' (3, 5). Jekk taqbel, jiena bdejt niċċaqlaq l{-}oġġetti tiegħi lejn dawn il{-}posizzjonijiet."} \\
\\ 
\texttt{X'inhu l{-}ewwel kmand tiegħek?} \\
            }
        }
    }
    & \\ \\

    \theutterance \stepcounter{utterance}  
    & & & \multicolumn{4}{p{0.6\linewidth}}{
        \cellcolor[rgb]{0.9,0.9,0.9}{
            \makecell[{{p{\linewidth}}}]{
                \texttt{\tiny{[GM$\langle$P2]}}
                \texttt{GĦID: Hello u grati għall{-}proposta! Ngħidilsek il{-}pożizzjonijiet attwali tiegħi qabel ma nibda nċaqlaq: C jinsab f' (1,5), L f' (3,2), u P f' (1,7). Ir{-}radici ta' C fil{-}bord tiegħi tippermetti li mmexxih sal{-}kaxxa t'isfel li tikkonfina miegħu, jiġifieri (1,6), imma L u P jeħtieġu miżuri aktar kumplessi biex jaslu fil{-}miri proposti tiegħek. Proponi li nkompli niddiskutu l{-}aqwa passaġġi possibbli, u jekk taqbel nibda niċċaqlaq C lejn (1,6) bħala l{-}ewwel pass tiegħi.} \\
            }
        }
    }
     \\ \\

    \theutterance \stepcounter{utterance}  
    & & & \multicolumn{2}{p{0.3\linewidth}}{
        \cellcolor[rgb]{0.9,0.9,0.9}{
            \makecell[{{p{\linewidth}}}]{
                \texttt{\tiny{[GM$|$GM]}}
                \texttt{GĦID: Hello u grati għall{-}proposta! Ngħidilsek il{-}pożizzjonijiet attwali tiegħi qabel ma nibda nċaqlaq: C jinsab f' (1,5), L f' (3,2), u P f' (1,7). Ir{-}radici ta' C fil{-}bord tiegħi tippermetti li mmexxih sal{-}kaxxa t'isfel li tikkonfina miegħu, jiġifieri (1,6), imma L u P jeħtieġu miżuri aktar kumplessi biex jaslu fil{-}miri proposti tiegħek. Proponi li nkompli niddiskutu l{-}aqwa passaġġi possibbli, u jekk taqbel nibda niċċaqlaq C lejn (1,6) bħala l{-}ewwel pass tiegħi.} \\
            }
        }
    }
    & & \\ \\

    \theutterance \stepcounter{utterance}  
    & & \multicolumn{4}{p{0.6\linewidth}}{
        \cellcolor[rgb]{0.9,0.9,0.9}{
            \makecell[{{p{\linewidth}}}]{
                \texttt{\tiny{[P1$\langle$GM]}}
                \texttt{Il{-}messaġġ tiegħek ġie mgħoddi lill{-}plejer l{-}ieħor.} \\
\\ 
\\ 
\texttt{Bħalissa qed tilagħbu r{-}round 1 minn massimu ta' 12.} \\
\texttt{Intom flimkien akkumulajtu 0 minn 8 penali.} \\
\texttt{Il{-}plejer l{-}ieħor bagħat dan il{-}messaġġ:} \\
\\ 
\\ 
\texttt{"Hello u grati għall{-}proposta! Ngħidilsek il{-}pożizzjonijiet attwali tiegħi qabel ma nibda nċaqlaq: C jinsab f' (1,5), L f' (3,2), u P f' (1,7). Ir{-}radici ta' C fil{-}bord tiegħi tippermetti li mmexxih sal{-}kaxxa t'isfel li tikkonfina miegħu, jiġifieri (1,6), imma L u P jeħtieġu miżuri aktar kumplessi biex jaslu fil{-}miri proposti tiegħek. Proponi li nkompli niddiskutu l{-}aqwa passaġġi possibbli, u jekk taqbel nibda niċċaqlaq C lejn (1,6) bħala l{-}ewwel pass tiegħi."} \\
\\ 
\\ 
\texttt{X'inhu l{-}kmand li jmiss tiegħek?} \\
            }
        }
    }
    & \\ \\

    \theutterance \stepcounter{utterance}  
    & \multicolumn{4}{p{0.6\linewidth}}{
        \cellcolor[rgb]{0.9,0.9,0.9}{
            \makecell[{{p{\linewidth}}}]{
                \texttt{\tiny{[P1$\rangle$GM]}}
                \texttt{GĦID: Nixtieq naqbel mal{-}ewwel pass tiegħek biex iċ{-}ċaqliq tiegħek jikkonċernja biss reġjuni aċċessibbli. Nipproponi li noqogħdu għal Konfigurazzjoni Komuni (CK) sempliċi ġol{-}futtur: P fil{-}bażi tax{-}xellug, C fil{-}qiegħ, u L fil{-}lemin estrem. Jiena se nibda billi nċaqlaq il{-}C tiegħi lejn (1,3) u r{-}reatt tiegħek lejn (1,3) kif mressqa.} \\
            }
        }
    }
    & & \\ \\

    \theutterance \stepcounter{utterance}  
    & & & \multicolumn{2}{p{0.3\linewidth}}{
        \cellcolor[rgb]{0.9,0.9,0.9}{
            \makecell[{{p{\linewidth}}}]{
                \texttt{\tiny{[GM$|$GM]}}
                \texttt{GĦID: Nixtieq naqbel mal{-}ewwel pass tiegħek biex iċ{-}ċaqliq tiegħek jikkonċernja biss reġjuni aċċessibbli. Nipproponi li noqogħdu għal Konfigurazzjoni Komuni (CK) sempliċi ġol{-}futtur: P fil{-}bażi tax{-}xellug, C fil{-}qiegħ, u L fil{-}lemin estrem. Jiena se nibda billi nċaqlaq il{-}C tiegħi lejn (1,3) u r{-}reatt tiegħek lejn (1,3) kif mressqa.} \\
            }
        }
    }
    & & \\ \\

    \theutterance \stepcounter{utterance}  
    & & \multicolumn{4}{p{0.6\linewidth}}{
        \cellcolor[rgb]{0.9,0.9,0.9}{
            \makecell[{{p{\linewidth}}}]{
                \texttt{\tiny{[GM$\rangle$P2]}}
                \texttt{Il{-}messaġġ tiegħek ġie mgħoddi lill{-}plejer l{-}ieħor.} \\
\\ 
\\ 
\texttt{Bħalissa qed tilagħbu r{-}round 2 minn massimu ta' 12.} \\
\texttt{Intom flimkien akkumulajtu 0 minn 8 penali.} \\
\texttt{Il{-}plejer l{-}ieħor bagħat dan il{-}messaġġ:} \\
\\ 
\\ 
\texttt{"Nixtieq naqbel mal{-}ewwel pass tiegħek biex iċ{-}ċaqliq tiegħek jikkonċernja biss reġjuni aċċessibbli. Nipproponi li noqogħdu għal Konfigurazzjoni Komuni (CK) sempliċi ġol{-}futtur: P fil{-}bażi tax{-}xellug, C fil{-}qiegħ, u L fil{-}lemin estrem. Jiena se nibda billi nċaqlaq il{-}C tiegħi lejn (1,3) u r{-}reatt tiegħek lejn (1,3) kif mressqa."} \\
\\ 
\\ 
\texttt{X'inhu l{-}kmand li jmiss tiegħek?} \\
            }
        }
    }
    & \\ \\

    \theutterance \stepcounter{utterance}  
    & & & \multicolumn{4}{p{0.6\linewidth}}{
        \cellcolor[rgb]{0.9,0.9,0.9}{
            \makecell[{{p{\linewidth}}}]{
                \texttt{\tiny{[GM$\langle$P2]}}
                \texttt{ĊAQLEQ: C, (1, 6)} \\
            }
        }
    }
     \\ \\

    \theutterance \stepcounter{utterance}  
    & & & \multicolumn{2}{p{0.3\linewidth}}{
        \cellcolor[rgb]{0.9,0.9,0.9}{
            \makecell[{{p{\linewidth}}}]{
                \texttt{\tiny{[GM$|$GM]}}
                \texttt{ĊAQLEQ: C, (1, 6)} \\
            }
        }
    }
    & & \\ \\

    \theutterance \stepcounter{utterance}  
    & & & \multicolumn{2}{p{0.3\linewidth}}{
        \cellcolor[rgb]{0.9,0.9,0.9}{
            \makecell[{{p{\linewidth}}}]{
                \texttt{\tiny{[GM$|$GM]}}
                \texttt{Penali: (1,6) mhijiex vojta, iżda fiha '─'.} \\
\\ 
\texttt{Intom flimkien akkumulajtu 1 minn 8 penali.} \\
\texttt{Jekk jogħġbok erġa' pprova!} \\
            }
        }
    }
    & & \\ \\

    \theutterance \stepcounter{utterance}  
    & & \multicolumn{4}{p{0.6\linewidth}}{
        \cellcolor[rgb]{0.9,0.9,0.9}{
            \makecell[{{p{\linewidth}}}]{
                \texttt{\tiny{[GM$\rangle$P2]}}
                \texttt{Penali: (1,6) mhijiex vojta, iżda fiha '─'.} \\
\\ 
\texttt{Intom flimkien akkumulajtu 1 minn 8 penali.} \\
\texttt{Jekk jogħġbok erġa' pprova!} \\
            }
        }
    }
    & \\ \\

    \theutterance \stepcounter{utterance}  
    & & & \multicolumn{4}{p{0.6\linewidth}}{
        \cellcolor[rgb]{0.9,0.9,0.9}{
            \makecell[{{p{\linewidth}}}]{
                \texttt{\tiny{[GM$\langle$P2]}}
                \texttt{GĦID: Jiddispjaċini għall{-}penali! Ma kellix inqis li l{-}kaxxa (1,6) fiha linja '─' mhux espressa fil{-}mappa viżiva. Nerġa' nieħu ħsieb il{-}passaġġi validi tiegħi. Il{-}bord tiegħi jidher bħali dan:} \\
\\ 
\texttt{```} \\
\texttt{ 1234567} \\
\texttt{╔═══╤═╤═╗} \\
\texttt{║◌◌◌│◌│◌║ 1} \\
\texttt{║◌◌L└─┤◌║ 2} \\
\texttt{║◌◌◌◌◌├─╢ 3} \\
\texttt{╟─┐◌┌─┤◌║ 4} \\
\texttt{║C│◌│◌│◌║ 5} \\
\texttt{╟─┼─┼─┼─╢ 6} \\
\texttt{║P│◌│◌│◌║ 7} \\
\texttt{╚═╧═╧═╧═╝} \\
\texttt{```} \\
\\ 
\texttt{P is (1,7), C is (1,5), and L is (3,2). Jidher li l{-}kaxxi ġo d{-}dawk iffighter huma l{-}uniki vojta. Nipproponi CK sempliċi li nistgħu nwasslu għaliha: C f' (2,5), L f' (3,5), u P f' (4,7). Jiena se nibda niċċaqlaq C tiegħi għal (2,5). Tista' tikkonferma li C u P huma fil{-}baxx fuq ix{-}xellug tiegħek, b'L fil{-}lemin estrem?} \\
            }
        }
    }
     \\ \\

    \theutterance \stepcounter{utterance}  
    & & & \multicolumn{2}{p{0.3\linewidth}}{
        \cellcolor[rgb]{0.9,0.9,0.9}{
            \makecell[{{p{\linewidth}}}]{
                \texttt{\tiny{[GM$|$GM]}}
                \texttt{GĦID: Jiddispjaċini għall{-}penali! Ma kellix inqis li l{-}kaxxa (1,6) fiha linja '─' mhux espressa fil{-}mappa viżiva. Nerġa' nieħu ħsieb il{-}passaġġi validi tiegħi. Il{-}bord tiegħi jidher bħali dan:} \\
\\ 
\texttt{```} \\
\texttt{ 1234567} \\
\texttt{╔═══╤═╤═╗} \\
\texttt{║◌◌◌│◌│◌║ 1} \\
\texttt{║◌◌L└─┤◌║ 2} \\
\texttt{║◌◌◌◌◌├─╢ 3} \\
\texttt{╟─┐◌┌─┤◌║ 4} \\
\texttt{║C│◌│◌│◌║ 5} \\
\texttt{╟─┼─┼─┼─╢ 6} \\
\texttt{║P│◌│◌│◌║ 7} \\
\texttt{╚═╧═╧═╧═╝} \\
\texttt{```} \\
\\ 
\texttt{P is (1,7), C is (1,5), and L is (3,2). Jidher li l{-}kaxxi ġo d{-}dawk iffighter huma l{-}uniki vojta. Nipproponi CK sempliċi li nistgħu nwasslu għaliha: C f' (2,5), L f' (3,5), u P f' (4,7). Jiena se nibda niċċaqlaq C tiegħi għal (2,5). Tista' tikkonferma li C u P huma fil{-}baxx fuq ix{-}xellug tiegħek, b'L fil{-}lemin estrem?} \\
            }
        }
    }
    & & \\ \\

    \theutterance \stepcounter{utterance}  
    & & & \multicolumn{2}{p{0.3\linewidth}}{
        \cellcolor[rgb]{0.9,0.9,0.9}{
            \makecell[{{p{\linewidth}}}]{
                \texttt{\tiny{[GM$|$GM]}}
                \texttt{Invalid format: tail is not empty: '} \\
\\ 
\\ 
\texttt{ 1234567} \\
\texttt{╔═══╤═╤═╗} \\
\texttt{║◌◌◌│◌│◌║ 1} \\
\texttt{║◌◌L└─┤◌║ 2} \\
\texttt{║◌◌◌◌◌├─╢ 3} \\
\texttt{╟─┐◌┌─┤◌║ 4} \\
\texttt{║C│◌│◌│◌║ 5} \\
\texttt{╟─┼─┼─┼─╢ 6} \\
\texttt{║P│◌│◌│◌║ 7} \\
\texttt{╚═╧═╧═╧═╝} \\
\\ 
\\ 
\texttt{P is (1,7), C is (1,5), and L is (3,2). Jidher li l{-}kaxxi ġo d{-}dawk iffighter huma l{-}uniki vojta. Nipproponi CK sempliċi li nistgħu nwasslu għaliha: C f' (2,5), L f' (3,5), u P f' (4,7). Jiena se nibda niċċaqlaq C tiegħi għal (2,5). Tista' tikkonferma li C u P huma fil{-}baxx fuq ix{-}xellug tiegħek, b'L fil{-}lemin estrem?} \\
            }
        }
    }
    & & \\ \\

    \theutterance \stepcounter{utterance}  
    & & \multicolumn{4}{p{0.6\linewidth}}{
        \cellcolor[rgb]{0.9,0.9,0.9}{
            \makecell[{{p{\linewidth}}}]{
                \texttt{\tiny{[GM$\rangle$P2]}}
                \texttt{Penali: Il{-}messaġġ tiegħek m'għandux ikun fih xejn wara l{-}kmand!} \\
\\ 
\texttt{Kun żgur li r{-}risposta tiegħek fiha biss jew `GĦID: <MESSAĠĠ>` jew `ĊAQLEQ: <OĠĠETT>, (<X>, <Y>)`, u xejn ieħor!} \\
\\ 
\texttt{Intom flimkien akkumulajtu 2 minn 8 penali.} \\
\texttt{Jekk jogħġbok erġa' pprova!} \\
            }
        }
    }
    & \\ \\

    \theutterance \stepcounter{utterance}  
    & & & \multicolumn{4}{p{0.6\linewidth}}{
        \cellcolor[rgb]{0.9,0.9,0.9}{
            \makecell[{{p{\linewidth}}}]{
                \texttt{\tiny{[GM$\langle$P2]}}
                \texttt{ĊAQLEQ: C, (2, 5)} \\
            }
        }
    }
     \\ \\

    \theutterance \stepcounter{utterance}  
    & & & \multicolumn{2}{p{0.3\linewidth}}{
        \cellcolor[rgb]{0.9,0.9,0.9}{
            \makecell[{{p{\linewidth}}}]{
                \texttt{\tiny{[GM$|$GM]}}
                \texttt{ĊAQLEQ: C, (2, 5)} \\
            }
        }
    }
    & & \\ \\

    \theutterance \stepcounter{utterance}  
    & & & \multicolumn{2}{p{0.3\linewidth}}{
        \cellcolor[rgb]{0.9,0.9,0.9}{
            \makecell[{{p{\linewidth}}}]{
                \texttt{\tiny{[GM$|$GM]}}
                \texttt{Penali: (2,5) mhijiex vojta, iżda fiha '│'.} \\
\\ 
\texttt{Intom flimkien akkumulajtu 3 minn 8 penali.} \\
\texttt{Jekk jogħġbok erġa' pprova!} \\
            }
        }
    }
    & & \\ \\

    \theutterance \stepcounter{utterance}  
    & & \multicolumn{4}{p{0.6\linewidth}}{
        \cellcolor[rgb]{0.9,0.9,0.9}{
            \makecell[{{p{\linewidth}}}]{
                \texttt{\tiny{[GM$\rangle$P2]}}
                \texttt{Penali: (2,5) mhijiex vojta, iżda fiha '│'.} \\
\\ 
\texttt{Intom flimkien akkumulajtu 3 minn 8 penali.} \\
\texttt{Jekk jogħġbok erġa' pprova!} \\
            }
        }
    }
    & \\ \\

    \theutterance \stepcounter{utterance}  
    & & & \multicolumn{4}{p{0.6\linewidth}}{
        \cellcolor[rgb]{0.9,0.9,0.9}{
            \makecell[{{p{\linewidth}}}]{
                \texttt{\tiny{[GM$\langle$P2]}}
                \texttt{GĦID: Jiddispjaċini għall{-}iżbalji. F'dan il{-}bord, l{-}unika ċelloli li jidhru kompletament vojta huma dawk bil{-}punt '◌'. Il{-}koordinati tagħhom huma: (1,1), (2,1), (3,1), (5,1), (7,1), (3,4), (3,5), (3,6), (4,5), (4,6), (4,7), (5,7), (6,7), u (7,7). Nipproponi l{-}Konfigurazzjoni Komuni segwenti: C f' (3,4), L f' (4,5), u P f' (4,7). Jiena immur iċċaqlaq il{-}C tiegħi għal (3,4) issa.} \\
            }
        }
    }
     \\ \\

    \theutterance \stepcounter{utterance}  
    & & & \multicolumn{2}{p{0.3\linewidth}}{
        \cellcolor[rgb]{0.9,0.9,0.9}{
            \makecell[{{p{\linewidth}}}]{
                \texttt{\tiny{[GM$|$GM]}}
                \texttt{GĦID: Jiddispjaċini għall{-}iżbalji. F'dan il{-}bord, l{-}unika ċelloli li jidhru kompletament vojta huma dawk bil{-}punt '◌'. Il{-}koordinati tagħhom huma: (1,1), (2,1), (3,1), (5,1), (7,1), (3,4), (3,5), (3,6), (4,5), (4,6), (4,7), (5,7), (6,7), u (7,7). Nipproponi l{-}Konfigurazzjoni Komuni segwenti: C f' (3,4), L f' (4,5), u P f' (4,7). Jiena immur iċċaqlaq il{-}C tiegħi għal (3,4) issa.} \\
            }
        }
    }
    & & \\ \\

    \theutterance \stepcounter{utterance}  
    & & \multicolumn{4}{p{0.6\linewidth}}{
        \cellcolor[rgb]{0.9,0.9,0.9}{
            \makecell[{{p{\linewidth}}}]{
                \texttt{\tiny{[P1$\langle$GM]}}
                \texttt{Il{-}messaġġ tiegħek ġie mgħoddi lill{-}plejer l{-}ieħor.} \\
\\ 
\\ 
\texttt{Bħalissa qed tilagħbu r{-}round 2 minn massimu ta' 12.} \\
\texttt{Intom flimkien akkumulajtu 3 minn 8 penali.} \\
\texttt{Il{-}plejer l{-}ieħor bagħat dan il{-}messaġġ:} \\
\\ 
\\ 
\texttt{"Jiddispjaċini għall{-}iżbalji. F'dan il{-}bord, l{-}unika ċelloli li jidhru kompletament vojta huma dawk bil{-}punt '◌'. Il{-}koordinati tagħhom huma: (1,1), (2,1), (3,1), (5,1), (7,1), (3,4), (3,5), (3,6), (4,5), (4,6), (4,7), (5,7), (6,7), u (7,7). Nipproponi l{-}Konfigurazzjoni Komuni segwenti: C f' (3,4), L f' (4,5), u P f' (4,7). Jiena immur iċċaqlaq il{-}C tiegħi għal (3,4) issa."} \\
\\ 
\\ 
\texttt{X'inhu l{-}kmand li jmiss tiegħek?} \\
            }
        }
    }
    & \\ \\

    \theutterance \stepcounter{utterance}  
    & \multicolumn{4}{p{0.6\linewidth}}{
        \cellcolor[rgb]{0.9,0.9,0.9}{
            \makecell[{{p{\linewidth}}}]{
                \texttt{\tiny{[P1$\rangle$GM]}}
                \texttt{Kmand: ĊAQLEQ: C, (3,4)} \\
            }
        }
    }
    & & \\ \\

    \theutterance \stepcounter{utterance}  
    & & & \multicolumn{2}{p{0.3\linewidth}}{
        \cellcolor[rgb]{0.9,0.9,0.9}{
            \makecell[{{p{\linewidth}}}]{
                \texttt{\tiny{[GM$|$GM]}}
                \texttt{Kmand: ĊAQLEQ: C, (3,4)} \\
            }
        }
    }
    & & \\ \\

    \theutterance \stepcounter{utterance}  
    & & & \multicolumn{2}{p{0.3\linewidth}}{
        \cellcolor[rgb]{0.9,0.9,0.9}{
            \makecell[{{p{\linewidth}}}]{
                \texttt{\tiny{[GM$|$GM]}}
                \texttt{Invalid format: head is not empty: 'Kmand: '} \\
            }
        }
    }
    & & \\ \\

    \theutterance \stepcounter{utterance}  
    & & \multicolumn{4}{p{0.6\linewidth}}{
        \cellcolor[rgb]{0.9,0.9,0.9}{
            \makecell[{{p{\linewidth}}}]{
                \texttt{\tiny{[P1$\langle$GM]}}
                \texttt{Penali: Il{-}messaġġ tiegħek m'għandux ikun fih xejn qabel il{-}kmand!} \\
\\ 
\texttt{Kun żgur li r{-}risposta tiegħek fiha biss jew `GĦID: <MESSAĠĠ>` jew `ĊAQLEQ: <OĠĠETT>, (<X>, <Y>)`, u xejn ieħor!} \\
\\ 
\texttt{Intom flimkien akkumulajtu 4 minn 8 penali.} \\
\texttt{Jekk jogħġbok erġa' pprova!} \\
            }
        }
    }
    & \\ \\

    \theutterance \stepcounter{utterance}  
    & \multicolumn{4}{p{0.6\linewidth}}{
        \cellcolor[rgb]{0.9,0.9,0.9}{
            \makecell[{{p{\linewidth}}}]{
                \texttt{\tiny{[P1$\rangle$GM]}}
                \texttt{ĊAQLEQ: C, (3,4)} \\
            }
        }
    }
    & & \\ \\

    \theutterance \stepcounter{utterance}  
    & & & \multicolumn{2}{p{0.3\linewidth}}{
        \cellcolor[rgb]{0.9,0.9,0.9}{
            \makecell[{{p{\linewidth}}}]{
                \texttt{\tiny{[GM$|$GM]}}
                \texttt{ĊAQLEQ: C, (3,4)} \\
            }
        }
    }
    & & \\ \\

    \theutterance \stepcounter{utterance}  
    & & & \multicolumn{2}{p{0.3\linewidth}}{
        \cellcolor[rgb]{0.9,0.9,0.9}{
            \makecell[{{p{\linewidth}}}]{
                \texttt{\tiny{[GM$|$GM]}}
                \texttt{C moved to (3, 4)} \\
            }
        }
    }
    & & \\ \\

    \theutterance \stepcounter{utterance}  
    & & \multicolumn{4}{p{0.6\linewidth}}{
        \cellcolor[rgb]{0.9,0.9,0.9}{
            \makecell[{{p{\linewidth}}}]{
                \texttt{\tiny{[GM$\rangle$P2]}}
                \texttt{Il{-}messaġġ tiegħek ġie mgħoddi lill{-}plejer l{-}ieħor.} \\
\\ 
\\ 
\texttt{Bħalissa qed tilagħbu r{-}round 3 minn massimu ta' 12.} \\
\texttt{Intom flimkien akkumulajtu 4 minn 8 penali.} \\
\texttt{Il{-}plejer l{-}ieħor mexxa oġġett fuq il{-}grilja tiegħu.} \\
\\ 
\\ 
\texttt{X'inhu l{-}kmand li jmiss tiegħek?} \\
            }
        }
    }
    & \\ \\

    \theutterance \stepcounter{utterance}  
    & & & \multicolumn{4}{p{0.6\linewidth}}{
        \cellcolor[rgb]{0.9,0.9,0.9}{
            \makecell[{{p{\linewidth}}}]{
                \texttt{\tiny{[GM$\langle$P2]}}
                \texttt{ĊAQLEQ: C, (3, 4)} \\
            }
        }
    }
     \\ \\

    \theutterance \stepcounter{utterance}  
    & & & \multicolumn{2}{p{0.3\linewidth}}{
        \cellcolor[rgb]{0.9,0.9,0.9}{
            \makecell[{{p{\linewidth}}}]{
                \texttt{\tiny{[GM$|$GM]}}
                \texttt{ĊAQLEQ: C, (3, 4)} \\
            }
        }
    }
    & & \\ \\

    \theutterance \stepcounter{utterance}  
    & & & \multicolumn{2}{p{0.3\linewidth}}{
        \cellcolor[rgb]{0.9,0.9,0.9}{
            \makecell[{{p{\linewidth}}}]{
                \texttt{\tiny{[GM$|$GM]}}
                \texttt{C moved to (3, 4)} \\
            }
        }
    }
    & & \\ \\

    \theutterance \stepcounter{utterance}  
    & & \multicolumn{4}{p{0.6\linewidth}}{
        \cellcolor[rgb]{0.9,0.9,0.9}{
            \makecell[{{p{\linewidth}}}]{
                \texttt{\tiny{[P1$\langle$GM]}}
                \texttt{Mexxejt 'C' għal (3,4) b'suċċess. Il{-}grilja aġġornata tiegħek tidher hekk:} \\
\texttt{```} \\
\texttt{ 1234567} \\
\texttt{╔═══╤═╤═╗} \\
\texttt{║◌◌◌│◌│◌║ 1} \\
\texttt{║◌P◌└─┤◌║ 2} \\
\texttt{║◌◌◌◌◌├─╢ 3} \\
\texttt{╟─┐C┌─┤◌║ 4} \\
\texttt{║◌│◌│◌│◌║ 5} \\
\texttt{╟─┼─┼─┼─╢ 6} \\
\texttt{║◌│◌│◌│L║ 7} \\
\texttt{╚═╧═╧═╧═╝} \\
\texttt{```} \\
\texttt{Bħalissa qed tilagħbu r{-}round 3 minn massimu ta' 12.} \\
\texttt{Intom flimkien akkumulajtu 4 minn 8 penali.} \\
\texttt{Il{-}plejer l{-}ieħor mexxa oġġett fuq il{-}grilja tiegħu.} \\
\\ 
\\ 
\texttt{X'inhu l{-}kmand li jmiss tiegħek?} \\
            }
        }
    }
    & \\ \\

    \theutterance \stepcounter{utterance}  
    & \multicolumn{4}{p{0.6\linewidth}}{
        \cellcolor[rgb]{0.9,0.9,0.9}{
            \makecell[{{p{\linewidth}}}]{
                \texttt{\tiny{[P1$\rangle$GM]}}
                \texttt{GĦID: Suċċess! C tiegħi huwa f' (3,4) wkoll. Irrenjew l{-}aqwa waqtien pass tiegħek. L{-}għan huwa Issa trid tlaqqaq P lejn (4,7) u L lejn (4,5). Imma peress li l{-}iskoperta taċ{-}ċelloli viċini mhix ċert, jiena se nagħmel ċaqlieq responsabbli. L{-}ewwel, se mbagħad nieħu ċ{-}ċaqlieq tiegħi għal P lejn (3,5). Jista' jkun li jkollna bżonn nirranġaw l{-}għan wara l{-}ewwel pass.} \\
            }
        }
    }
    & & \\ \\

    \theutterance \stepcounter{utterance}  
    & & & \multicolumn{2}{p{0.3\linewidth}}{
        \cellcolor[rgb]{0.9,0.9,0.9}{
            \makecell[{{p{\linewidth}}}]{
                \texttt{\tiny{[GM$|$GM]}}
                \texttt{GĦID: Suċċess! C tiegħi huwa f' (3,4) wkoll. Irrenjew l{-}aqwa waqtien pass tiegħek. L{-}għan huwa Issa trid tlaqqaq P lejn (4,7) u L lejn (4,5). Imma peress li l{-}iskoperta taċ{-}ċelloli viċini mhix ċert, jiena se nagħmel ċaqlieq responsabbli. L{-}ewwel, se mbagħad nieħu ċ{-}ċaqlieq tiegħi għal P lejn (3,5). Jista' jkun li jkollna bżonn nirranġaw l{-}għan wara l{-}ewwel pass.} \\
            }
        }
    }
    & & \\ \\

    \theutterance \stepcounter{utterance}  
    & & \multicolumn{4}{p{0.6\linewidth}}{
        \cellcolor[rgb]{0.9,0.9,0.9}{
            \makecell[{{p{\linewidth}}}]{
                \texttt{\tiny{[GM$\rangle$P2]}}
                \texttt{Mexxejt 'C' għal (3,4) b'suċċess. Il{-}grilja aġġornata tiegħek tidher hekk:} \\
\texttt{```} \\
\texttt{ 1234567} \\
\texttt{╔═══╤═╤═╗} \\
\texttt{║◌◌◌│◌│◌║ 1} \\
\texttt{║◌◌L└─┤◌║ 2} \\
\texttt{║◌◌◌◌◌├─╢ 3} \\
\texttt{╟─┐C┌─┤◌║ 4} \\
\texttt{║◌│◌│◌│◌║ 5} \\
\texttt{╟─┼─┼─┼─╢ 6} \\
\texttt{║P│◌│◌│◌║ 7} \\
\texttt{╚═╧═╧═╧═╝} \\
\texttt{```} \\
\texttt{Bħalissa qed tilagħbu r{-}round 4 minn massimu ta' 12.} \\
\texttt{Intom flimkien akkumulajtu 4 minn 8 penali.} \\
\texttt{Il{-}plejer l{-}ieħor bagħat dan il{-}messaġġ:} \\
\\ 
\\ 
\texttt{"Suċċess! C tiegħi huwa f' (3,4) wkoll. Irrenjew l{-}aqwa waqtien pass tiegħek. L{-}għan huwa Issa trid tlaqqaq P lejn (4,7) u L lejn (4,5). Imma peress li l{-}iskoperta taċ{-}ċelloli viċini mhix ċert, jiena se nagħmel ċaqlieq responsabbli. L{-}ewwel, se mbagħad nieħu ċ{-}ċaqlieq tiegħi għal P lejn (3,5). Jista' jkun li jkollna bżonn nirranġaw l{-}għan wara l{-}ewwel pass."} \\
\\ 
\\ 
\texttt{X'inhu l{-}kmand li jmiss tiegħek?} \\
            }
        }
    }
    & \\ \\

    \theutterance \stepcounter{utterance}  
    & & & \multicolumn{4}{p{0.6\linewidth}}{
        \cellcolor[rgb]{0.9,0.9,0.9}{
            \makecell[{{p{\linewidth}}}]{
                \texttt{\tiny{[GM$\langle$P2]}}
                \texttt{GĦID: Prosit għall{-}ftehim fuq C! Se nimxi mat{-}tweġiba tiegħek: P għal (3,5) u L għal (4,5) huma miżuri tajbin. L{-}ewwel se niċċaqlaq P tiegħi lejn (3,5).} \\
            }
        }
    }
     \\ \\

    \theutterance \stepcounter{utterance}  
    & & & \multicolumn{2}{p{0.3\linewidth}}{
        \cellcolor[rgb]{0.9,0.9,0.9}{
            \makecell[{{p{\linewidth}}}]{
                \texttt{\tiny{[GM$|$GM]}}
                \texttt{GĦID: Prosit għall{-}ftehim fuq C! Se nimxi mat{-}tweġiba tiegħek: P għal (3,5) u L għal (4,5) huma miżuri tajbin. L{-}ewwel se niċċaqlaq P tiegħi lejn (3,5).} \\
            }
        }
    }
    & & \\ \\

    \theutterance \stepcounter{utterance}  
    & & \multicolumn{4}{p{0.6\linewidth}}{
        \cellcolor[rgb]{0.9,0.9,0.9}{
            \makecell[{{p{\linewidth}}}]{
                \texttt{\tiny{[P1$\langle$GM]}}
                \texttt{Il{-}messaġġ tiegħek ġie mgħoddi lill{-}plejer l{-}ieħor.} \\
\\ 
\\ 
\texttt{Bħalissa qed tilagħbu r{-}round 4 minn massimu ta' 12.} \\
\texttt{Intom flimkien akkumulajtu 4 minn 8 penali.} \\
\texttt{Il{-}plejer l{-}ieħor bagħat dan il{-}messaġġ:} \\
\\ 
\\ 
\texttt{"Prosit għall{-}ftehim fuq C! Se nimxi mat{-}tweġiba tiegħek: P għal (3,5) u L għal (4,5) huma miżuri tajbin. L{-}ewwel se niċċaqlaq P tiegħi lejn (3,5)."} \\
\\ 
\\ 
\texttt{X'inhu l{-}kmand li jmiss tiegħek?} \\
            }
        }
    }
    & \\ \\

    \theutterance \stepcounter{utterance}  
    & \multicolumn{4}{p{0.6\linewidth}}{
        \cellcolor[rgb]{0.9,0.9,0.9}{
            \makecell[{{p{\linewidth}}}]{
                \texttt{\tiny{[P1$\rangle$GM]}}
                \texttt{ĊAQLEQ: P, (3,5)} \\
            }
        }
    }
    & & \\ \\

    \theutterance \stepcounter{utterance}  
    & & & \multicolumn{2}{p{0.3\linewidth}}{
        \cellcolor[rgb]{0.9,0.9,0.9}{
            \makecell[{{p{\linewidth}}}]{
                \texttt{\tiny{[GM$|$GM]}}
                \texttt{ĊAQLEQ: P, (3,5)} \\
            }
        }
    }
    & & \\ \\

    \theutterance \stepcounter{utterance}  
    & & & \multicolumn{2}{p{0.3\linewidth}}{
        \cellcolor[rgb]{0.9,0.9,0.9}{
            \makecell[{{p{\linewidth}}}]{
                \texttt{\tiny{[GM$|$GM]}}
                \texttt{P moved to (3, 5)} \\
            }
        }
    }
    & & \\ \\

    \theutterance \stepcounter{utterance}  
    & & \multicolumn{4}{p{0.6\linewidth}}{
        \cellcolor[rgb]{0.9,0.9,0.9}{
            \makecell[{{p{\linewidth}}}]{
                \texttt{\tiny{[GM$\rangle$P2]}}
                \texttt{Il{-}messaġġ tiegħek ġie mgħoddi lill{-}plejer l{-}ieħor.} \\
\\ 
\\ 
\texttt{Bħalissa qed tilagħbu r{-}round 5 minn massimu ta' 12.} \\
\texttt{Intom flimkien akkumulajtu 4 minn 8 penali.} \\
\texttt{Il{-}plejer l{-}ieħor mexxa oġġett fuq il{-}grilja tiegħu.} \\
\\ 
\\ 
\texttt{X'inhu l{-}kmand li jmiss tiegħek?} \\
            }
        }
    }
    & \\ \\

    \theutterance \stepcounter{utterance}  
    & & & \multicolumn{4}{p{0.6\linewidth}}{
        \cellcolor[rgb]{0.9,0.9,0.9}{
            \makecell[{{p{\linewidth}}}]{
                \texttt{\tiny{[GM$\langle$P2]}}
                \texttt{ĊAQLEQ: P, (3, 5)} \\
            }
        }
    }
     \\ \\

    \theutterance \stepcounter{utterance}  
    & & & \multicolumn{2}{p{0.3\linewidth}}{
        \cellcolor[rgb]{0.9,0.9,0.9}{
            \makecell[{{p{\linewidth}}}]{
                \texttt{\tiny{[GM$|$GM]}}
                \texttt{ĊAQLEQ: P, (3, 5)} \\
            }
        }
    }
    & & \\ \\

    \theutterance \stepcounter{utterance}  
    & & & \multicolumn{2}{p{0.3\linewidth}}{
        \cellcolor[rgb]{0.9,0.9,0.9}{
            \makecell[{{p{\linewidth}}}]{
                \texttt{\tiny{[GM$|$GM]}}
                \texttt{P moved to (3, 5)} \\
            }
        }
    }
    & & \\ \\

    \theutterance \stepcounter{utterance}  
    & & \multicolumn{4}{p{0.6\linewidth}}{
        \cellcolor[rgb]{0.9,0.9,0.9}{
            \makecell[{{p{\linewidth}}}]{
                \texttt{\tiny{[P1$\langle$GM]}}
                \texttt{Mexxejt 'P' għal (3,5) b'suċċess. Il{-}grilja aġġornata tiegħek tidher hekk:} \\
\texttt{```} \\
\texttt{ 1234567} \\
\texttt{╔═══╤═╤═╗} \\
\texttt{║◌◌◌│◌│◌║ 1} \\
\texttt{║◌◌◌└─┤◌║ 2} \\
\texttt{║◌◌◌◌◌├─╢ 3} \\
\texttt{╟─┐C┌─┤◌║ 4} \\
\texttt{║◌│P│◌│◌║ 5} \\
\texttt{╟─┼─┼─┼─╢ 6} \\
\texttt{║◌│◌│◌│L║ 7} \\
\texttt{╚═╧═╧═╧═╝} \\
\texttt{```} \\
\texttt{Bħalissa qed tilagħbu r{-}round 5 minn massimu ta' 12.} \\
\texttt{Intom flimkien akkumulajtu 4 minn 8 penali.} \\
\texttt{Il{-}plejer l{-}ieħor mexxa oġġett fuq il{-}grilja tiegħu.} \\
\\ 
\\ 
\texttt{X'inhu l{-}kmand li jmiss tiegħek?} \\
            }
        }
    }
    & \\ \\

    \theutterance \stepcounter{utterance}  
    & \multicolumn{4}{p{0.6\linewidth}}{
        \cellcolor[rgb]{0.9,0.9,0.9}{
            \makecell[{{p{\linewidth}}}]{
                \texttt{\tiny{[P1$\rangle$GM]}}
                \texttt{GĦID: P huwa f' (3,5) b'suċċess! L{-}aħħar oġġett huwa L. B'din il{-}grilja, naħseb li l{-}aqwa għan għal L huwa (7,7) minħabba li hemm aċċess ovvju. Inti qbil? Jekk iva, jiena se niċċaqlaq L tiegħi għal (7,7) fir{-}round 6, u inti tagħmel l{-}istess biex inallinejawh.} \\
            }
        }
    }
    & & \\ \\

    \theutterance \stepcounter{utterance}  
    & & & \multicolumn{2}{p{0.3\linewidth}}{
        \cellcolor[rgb]{0.9,0.9,0.9}{
            \makecell[{{p{\linewidth}}}]{
                \texttt{\tiny{[GM$|$GM]}}
                \texttt{GĦID: P huwa f' (3,5) b'suċċess! L{-}aħħar oġġett huwa L. B'din il{-}grilja, naħseb li l{-}aqwa għan għal L huwa (7,7) minħabba li hemm aċċess ovvju. Inti qbil? Jekk iva, jiena se niċċaqlaq L tiegħi għal (7,7) fir{-}round 6, u inti tagħmel l{-}istess biex inallinejawh.} \\
            }
        }
    }
    & & \\ \\

    \theutterance \stepcounter{utterance}  
    & & \multicolumn{4}{p{0.6\linewidth}}{
        \cellcolor[rgb]{0.9,0.9,0.9}{
            \makecell[{{p{\linewidth}}}]{
                \texttt{\tiny{[GM$\rangle$P2]}}
                \texttt{Mexxejt 'P' għal (3,5) b'suċċess. Il{-}grilja aġġornata tiegħek tidher hekk:} \\
\texttt{```} \\
\texttt{ 1234567} \\
\texttt{╔═══╤═╤═╗} \\
\texttt{║◌◌◌│◌│◌║ 1} \\
\texttt{║◌◌L└─┤◌║ 2} \\
\texttt{║◌◌◌◌◌├─╢ 3} \\
\texttt{╟─┐C┌─┤◌║ 4} \\
\texttt{║◌│P│◌│◌║ 5} \\
\texttt{╟─┼─┼─┼─╢ 6} \\
\texttt{║◌│◌│◌│◌║ 7} \\
\texttt{╚═╧═╧═╧═╝} \\
\texttt{```} \\
\texttt{Bħalissa qed tilagħbu r{-}round 6 minn massimu ta' 12.} \\
\texttt{Intom flimkien akkumulajtu 4 minn 8 penali.} \\
\texttt{Il{-}plejer l{-}ieħor bagħat dan il{-}messaġġ:} \\
\\ 
\\ 
\texttt{"P huwa f' (3,5) b'suċċess! L{-}aħħar oġġett huwa L. B'din il{-}grilja, naħseb li l{-}aqwa għan għal L huwa (7,7) minħabba li hemm aċċess ovvju. Inti qbil? Jekk iva, jiena se niċċaqlaq L tiegħi għal (7,7) fir{-}round 6, u inti tagħmel l{-}istess biex inallinejawh."} \\
\\ 
\\ 
\texttt{X'inhu l{-}kmand li jmiss tiegħek?} \\
            }
        }
    }
    & \\ \\

    \theutterance \stepcounter{utterance}  
    & & & \multicolumn{4}{p{0.6\linewidth}}{
        \cellcolor[rgb]{0.9,0.9,0.9}{
            \makecell[{{p{\linewidth}}}]{
                \texttt{\tiny{[GM$\langle$P2]}}
                \texttt{GĦID: Ma nistax inqabel għal (7,7) għal L minħabba li L tiegħi jinsab imwaħħal f' (3,2) u ma jistax jiċċaqlaq direttament lejn (7,7). Jien se niċċaqlaq L għal (4,5) bħala l{-}għan finali tiegħi. Jista' jkun li tista' tpoġġi L tiegħek f' (4,5) wkoll?} \\
            }
        }
    }
     \\ \\

    \theutterance \stepcounter{utterance}  
    & & & \multicolumn{2}{p{0.3\linewidth}}{
        \cellcolor[rgb]{0.9,0.9,0.9}{
            \makecell[{{p{\linewidth}}}]{
                \texttt{\tiny{[GM$|$GM]}}
                \texttt{GĦID: Ma nistax inqabel għal (7,7) għal L minħabba li L tiegħi jinsab imwaħħal f' (3,2) u ma jistax jiċċaqlaq direttament lejn (7,7). Jien se niċċaqlaq L għal (4,5) bħala l{-}għan finali tiegħi. Jista' jkun li tista' tpoġġi L tiegħek f' (4,5) wkoll?} \\
            }
        }
    }
    & & \\ \\

    \theutterance \stepcounter{utterance}  
    & & \multicolumn{4}{p{0.6\linewidth}}{
        \cellcolor[rgb]{0.9,0.9,0.9}{
            \makecell[{{p{\linewidth}}}]{
                \texttt{\tiny{[P1$\langle$GM]}}
                \texttt{Il{-}messaġġ tiegħek ġie mgħoddi lill{-}plejer l{-}ieħor.} \\
\\ 
\\ 
\texttt{Bħalissa qed tilagħbu r{-}round 6 minn massimu ta' 12.} \\
\texttt{Intom flimkien akkumulajtu 4 minn 8 penali.} \\
\texttt{Il{-}plejer l{-}ieħor bagħat dan il{-}messaġġ:} \\
\\ 
\\ 
\texttt{"Ma nistax inqabel għal (7,7) għal L minħabba li L tiegħi jinsab imwaħħal f' (3,2) u ma jistax jiċċaqlaq direttament lejn (7,7). Jien se niċċaqlaq L għal (4,5) bħala l{-}għan finali tiegħi. Jista' jkun li tista' tpoġġi L tiegħek f' (4,5) wkoll?"} \\
\\ 
\\ 
\texttt{X'inhu l{-}kmand li jmiss tiegħek?} \\
            }
        }
    }
    & \\ \\

    \theutterance \stepcounter{utterance}  
    & \multicolumn{4}{p{0.6\linewidth}}{
        \cellcolor[rgb]{0.9,0.9,0.9}{
            \makecell[{{p{\linewidth}}}]{
                \texttt{\tiny{[P1$\rangle$GM]}}
                \texttt{GĦID: Jien naqbel, (4,5) huwa punt eċċellenti għall{-}għan finali tagħna. Nistennie ċ{-}ċaqliq tiegħek lejn (4,5) u jien ukoll se nagħmel l{-}istess biex niżguraw li L jallinja. Prosit għall{-}istrateġija!} \\
            }
        }
    }
    & & \\ \\

    \theutterance \stepcounter{utterance}  
    & & & \multicolumn{2}{p{0.3\linewidth}}{
        \cellcolor[rgb]{0.9,0.9,0.9}{
            \makecell[{{p{\linewidth}}}]{
                \texttt{\tiny{[GM$|$GM]}}
                \texttt{GĦID: Jien naqbel, (4,5) huwa punt eċċellenti għall{-}għan finali tagħna. Nistennie ċ{-}ċaqliq tiegħek lejn (4,5) u jien ukoll se nagħmel l{-}istess biex niżguraw li L jallinja. Prosit għall{-}istrateġija!} \\
            }
        }
    }
    & & \\ \\

    \theutterance \stepcounter{utterance}  
    & & \multicolumn{4}{p{0.6\linewidth}}{
        \cellcolor[rgb]{0.9,0.9,0.9}{
            \makecell[{{p{\linewidth}}}]{
                \texttt{\tiny{[GM$\rangle$P2]}}
                \texttt{Il{-}messaġġ tiegħek ġie mgħoddi lill{-}plejer l{-}ieħor.} \\
\\ 
\\ 
\texttt{Bħalissa qed tilagħbu r{-}round 7 minn massimu ta' 12.} \\
\texttt{Intom flimkien akkumulajtu 4 minn 8 penali.} \\
\texttt{Il{-}plejer l{-}ieħor bagħat dan il{-}messaġġ:} \\
\\ 
\\ 
\texttt{"Jien naqbel, (4,5) huwa punt eċċellenti għall{-}għan finali tagħna. Nistennie ċ{-}ċaqliq tiegħek lejn (4,5) u jien ukoll se nagħmel l{-}istess biex niżguraw li L jallinja. Prosit għall{-}istrateġija!"} \\
\\ 
\\ 
\texttt{X'inhu l{-}kmand li jmiss tiegħek?} \\
            }
        }
    }
    & \\ \\

    \theutterance \stepcounter{utterance}  
    & & & \multicolumn{4}{p{0.6\linewidth}}{
        \cellcolor[rgb]{0.9,0.9,0.9}{
            \makecell[{{p{\linewidth}}}]{
                \texttt{\tiny{[GM$\langle$P2]}}
                \texttt{ĊAQLEQ: L, (4, 5)} \\
            }
        }
    }
     \\ \\

    \theutterance \stepcounter{utterance}  
    & & & \multicolumn{2}{p{0.3\linewidth}}{
        \cellcolor[rgb]{0.9,0.9,0.9}{
            \makecell[{{p{\linewidth}}}]{
                \texttt{\tiny{[GM$|$GM]}}
                \texttt{ĊAQLEQ: L, (4, 5)} \\
            }
        }
    }
    & & \\ \\

    \theutterance \stepcounter{utterance}  
    & & & \multicolumn{2}{p{0.3\linewidth}}{
        \cellcolor[rgb]{0.9,0.9,0.9}{
            \makecell[{{p{\linewidth}}}]{
                \texttt{\tiny{[GM$|$GM]}}
                \texttt{Penali: (4,5) mhijiex vojta, iżda fiha '│'.} \\
\\ 
\texttt{Intom flimkien akkumulajtu 5 minn 8 penali.} \\
\texttt{Jekk jogħġbok erġa' pprova!} \\
            }
        }
    }
    & & \\ \\

    \theutterance \stepcounter{utterance}  
    & & \multicolumn{4}{p{0.6\linewidth}}{
        \cellcolor[rgb]{0.9,0.9,0.9}{
            \makecell[{{p{\linewidth}}}]{
                \texttt{\tiny{[GM$\rangle$P2]}}
                \texttt{Penali: (4,5) mhijiex vojta, iżda fiha '│'.} \\
\\ 
\texttt{Intom flimkien akkumulajtu 5 minn 8 penali.} \\
\texttt{Jekk jogħġbok erġa' pprova!} \\
            }
        }
    }
    & \\ \\

    \theutterance \stepcounter{utterance}  
    & & & \multicolumn{4}{p{0.6\linewidth}}{
        \cellcolor[rgb]{0.9,0.9,0.9}{
            \makecell[{{p{\linewidth}}}]{
                \texttt{\tiny{[GM$\langle$P2]}}
                \texttt{GĦID: Jiddispjaċini, laqtajt '│' f' (4,5). L{-}unika ċelloli vojta ħdejn L huma (1,1), (2,1), (3,1), (5,1), (7,1). Nipproponi L għal (3,1). Jista' jkun inti tpoġġi L tiegħek f' (3,1) ukoll?} \\
            }
        }
    }
     \\ \\

    \theutterance \stepcounter{utterance}  
    & & & \multicolumn{2}{p{0.3\linewidth}}{
        \cellcolor[rgb]{0.9,0.9,0.9}{
            \makecell[{{p{\linewidth}}}]{
                \texttt{\tiny{[GM$|$GM]}}
                \texttt{GĦID: Jiddispjaċini, laqtajt '│' f' (4,5). L{-}unika ċelloli vojta ħdejn L huma (1,1), (2,1), (3,1), (5,1), (7,1). Nipproponi L għal (3,1). Jista' jkun inti tpoġġi L tiegħek f' (3,1) ukoll?} \\
            }
        }
    }
    & & \\ \\

    \theutterance \stepcounter{utterance}  
    & & \multicolumn{4}{p{0.6\linewidth}}{
        \cellcolor[rgb]{0.9,0.9,0.9}{
            \makecell[{{p{\linewidth}}}]{
                \texttt{\tiny{[P1$\langle$GM]}}
                \texttt{Il{-}messaġġ tiegħek ġie mgħoddi lill{-}plejer l{-}ieħor.} \\
\\ 
\\ 
\texttt{Bħalissa qed tilagħbu r{-}round 7 minn massimu ta' 12.} \\
\texttt{Intom flimkien akkumulajtu 5 minn 8 penali.} \\
\texttt{Il{-}plejer l{-}ieħor bagħat dan il{-}messaġġ:} \\
\\ 
\\ 
\texttt{"Jiddispjaċini, laqtajt '│' f' (4,5). L{-}unika ċelloli vojta ħdejn L huma (1,1), (2,1), (3,1), (5,1), (7,1). Nipproponi L għal (3,1). Jista' jkun inti tpoġġi L tiegħek f' (3,1) ukoll?"} \\
\\ 
\\ 
\texttt{X'inhu l{-}kmand li jmiss tiegħek?} \\
            }
        }
    }
    & \\ \\

    \theutterance \stepcounter{utterance}  
    & \multicolumn{4}{p{0.6\linewidth}}{
        \cellcolor[rgb]{0.9,0.9,0.9}{
            \makecell[{{p{\linewidth}}}]{
                \texttt{\tiny{[P1$\rangle$GM]}}
                \texttt{GĦID: Niskatta liċ{-}ċellola (3,1) hija wkoll vojta għalija. Naqbel li npoġġu L f' (3,1). Jiena se niċċaqlaq L tiegħi minn (7,7) għal (3,1) biex nallinjaw ma' (3,1).} \\
            }
        }
    }
    & & \\ \\

    \theutterance \stepcounter{utterance}  
    & & & \multicolumn{2}{p{0.3\linewidth}}{
        \cellcolor[rgb]{0.9,0.9,0.9}{
            \makecell[{{p{\linewidth}}}]{
                \texttt{\tiny{[GM$|$GM]}}
                \texttt{GĦID: Niskatta liċ{-}ċellola (3,1) hija wkoll vojta għalija. Naqbel li npoġġu L f' (3,1). Jiena se niċċaqlaq L tiegħi minn (7,7) għal (3,1) biex nallinjaw ma' (3,1).} \\
            }
        }
    }
    & & \\ \\

    \theutterance \stepcounter{utterance}  
    & & \multicolumn{4}{p{0.6\linewidth}}{
        \cellcolor[rgb]{0.9,0.9,0.9}{
            \makecell[{{p{\linewidth}}}]{
                \texttt{\tiny{[GM$\rangle$P2]}}
                \texttt{Il{-}messaġġ tiegħek ġie mgħoddi lill{-}plejer l{-}ieħor.} \\
\\ 
\\ 
\texttt{Bħalissa qed tilagħbu r{-}round 8 minn massimu ta' 12.} \\
\texttt{Intom flimkien akkumulajtu 5 minn 8 penali.} \\
\texttt{Il{-}plejer l{-}ieħor bagħat dan il{-}messaġġ:} \\
\\ 
\\ 
\texttt{"Niskatta liċ{-}ċellola (3,1) hija wkoll vojta għalija. Naqbel li npoġġu L f' (3,1). Jiena se niċċaqlaq L tiegħi minn (7,7) għal (3,1) biex nallinjaw ma' (3,1)."} \\
\\ 
\\ 
\texttt{X'inhu l{-}kmand li jmiss tiegħek?} \\
            }
        }
    }
    & \\ \\

    \theutterance \stepcounter{utterance}  
    & & & \multicolumn{4}{p{0.6\linewidth}}{
        \cellcolor[rgb]{0.9,0.9,0.9}{
            \makecell[{{p{\linewidth}}}]{
                \texttt{\tiny{[GM$\langle$P2]}}
                \texttt{ĊAQLEQ: L, (3, 1)} \\
            }
        }
    }
     \\ \\

    \theutterance \stepcounter{utterance}  
    & & & \multicolumn{2}{p{0.3\linewidth}}{
        \cellcolor[rgb]{0.9,0.9,0.9}{
            \makecell[{{p{\linewidth}}}]{
                \texttt{\tiny{[GM$|$GM]}}
                \texttt{ĊAQLEQ: L, (3, 1)} \\
            }
        }
    }
    & & \\ \\

    \theutterance \stepcounter{utterance}  
    & & & \multicolumn{2}{p{0.3\linewidth}}{
        \cellcolor[rgb]{0.9,0.9,0.9}{
            \makecell[{{p{\linewidth}}}]{
                \texttt{\tiny{[GM$|$GM]}}
                \texttt{L moved to (3, 1)} \\
            }
        }
    }
    & & \\ \\

    \theutterance \stepcounter{utterance}  
    & & \multicolumn{4}{p{0.6\linewidth}}{
        \cellcolor[rgb]{0.9,0.9,0.9}{
            \makecell[{{p{\linewidth}}}]{
                \texttt{\tiny{[P1$\langle$GM]}}
                \texttt{Il{-}messaġġ tiegħek ġie mgħoddi lill{-}plejer l{-}ieħor.} \\
\\ 
\\ 
\texttt{Bħalissa qed tilagħbu r{-}round 8 minn massimu ta' 12.} \\
\texttt{Intom flimkien akkumulajtu 5 minn 8 penali.} \\
\texttt{Il{-}plejer l{-}ieħor mexxa oġġett fuq il{-}grilja tiegħu.} \\
\\ 
\\ 
\texttt{X'inhu l{-}kmand li jmiss tiegħek?} \\
            }
        }
    }
    & \\ \\

    \theutterance \stepcounter{utterance}  
    & \multicolumn{4}{p{0.6\linewidth}}{
        \cellcolor[rgb]{0.9,0.9,0.9}{
            \makecell[{{p{\linewidth}}}]{
                \texttt{\tiny{[P1$\rangle$GM]}}
                \texttt{ĊAQLEQ: L, (3,1)} \\
            }
        }
    }
    & & \\ \\

    \theutterance \stepcounter{utterance}  
    & & & \multicolumn{2}{p{0.3\linewidth}}{
        \cellcolor[rgb]{0.9,0.9,0.9}{
            \makecell[{{p{\linewidth}}}]{
                \texttt{\tiny{[GM$|$GM]}}
                \texttt{ĊAQLEQ: L, (3,1)} \\
            }
        }
    }
    & & \\ \\

    \theutterance \stepcounter{utterance}  
    & & & \multicolumn{2}{p{0.3\linewidth}}{
        \cellcolor[rgb]{0.9,0.9,0.9}{
            \makecell[{{p{\linewidth}}}]{
                \texttt{\tiny{[GM$|$GM]}}
                \texttt{L moved to (3, 1)} \\
            }
        }
    }
    & & \\ \\

    \theutterance \stepcounter{utterance}  
    & & \multicolumn{4}{p{0.6\linewidth}}{
        \cellcolor[rgb]{0.9,0.9,0.9}{
            \makecell[{{p{\linewidth}}}]{
                \texttt{\tiny{[GM$\rangle$P2]}}
                \texttt{Mexxejt 'L' għal (3,1) b'suċċess. Il{-}grilja aġġornata tiegħek tidher hekk:} \\
\texttt{```} \\
\texttt{ 1234567} \\
\texttt{╔═══╤═╤═╗} \\
\texttt{║◌◌L│◌│◌║ 1} \\
\texttt{║◌◌◌└─┤◌║ 2} \\
\texttt{║◌◌◌◌◌├─╢ 3} \\
\texttt{╟─┐C┌─┤◌║ 4} \\
\texttt{║◌│P│◌│◌║ 5} \\
\texttt{╟─┼─┼─┼─╢ 6} \\
\texttt{║◌│◌│◌│◌║ 7} \\
\texttt{╚═╧═╧═╧═╝} \\
\texttt{```} \\
\texttt{Bħalissa qed tilagħbu r{-}round 9 minn massimu ta' 12.} \\
\texttt{Intom flimkien akkumulajtu 5 minn 8 penali.} \\
\texttt{Il{-}plejer l{-}ieħor mexxa oġġett fuq il{-}grilja tiegħu.} \\
\\ 
\\ 
\texttt{X'inhu l{-}kmand li jmiss tiegħek?} \\
            }
        }
    }
    & \\ \\

    \theutterance \stepcounter{utterance}  
    & & & \multicolumn{4}{p{0.6\linewidth}}{
        \cellcolor[rgb]{0.9,0.9,0.9}{
            \makecell[{{p{\linewidth}}}]{
                \texttt{\tiny{[GM$\langle$P2]}}
                \texttt{GĦID: L huwa f' (3,1) b'suċċess! L{-}oġġetti kollha tiegħi huma mqiegħda: C f' (3,4), P f' (3,5), u L f' (3,1). Inti għandek L f' (3,1) wkoll? Jekk iva, lestejt?} \\
            }
        }
    }
     \\ \\

    \theutterance \stepcounter{utterance}  
    & & & \multicolumn{2}{p{0.3\linewidth}}{
        \cellcolor[rgb]{0.9,0.9,0.9}{
            \makecell[{{p{\linewidth}}}]{
                \texttt{\tiny{[GM$|$GM]}}
                \texttt{GĦID: L huwa f' (3,1) b'suċċess! L{-}oġġetti kollha tiegħi huma mqiegħda: C f' (3,4), P f' (3,5), u L f' (3,1). Inti għandek L f' (3,1) wkoll? Jekk iva, lestejt?} \\
            }
        }
    }
    & & \\ \\

    \theutterance \stepcounter{utterance}  
    & & \multicolumn{4}{p{0.6\linewidth}}{
        \cellcolor[rgb]{0.9,0.9,0.9}{
            \makecell[{{p{\linewidth}}}]{
                \texttt{\tiny{[P1$\langle$GM]}}
                \texttt{Mexxejt 'L' għal (3,1) b'suċċess. Il{-}grilja aġġornata tiegħek tidher hekk:} \\
\texttt{```} \\
\texttt{ 1234567} \\
\texttt{╔═══╤═╤═╗} \\
\texttt{║◌◌L│◌│◌║ 1} \\
\texttt{║◌◌◌└─┤◌║ 2} \\
\texttt{║◌◌◌◌◌├─╢ 3} \\
\texttt{╟─┐C┌─┤◌║ 4} \\
\texttt{║◌│P│◌│◌║ 5} \\
\texttt{╟─┼─┼─┼─╢ 6} \\
\texttt{║◌│◌│◌│◌║ 7} \\
\texttt{╚═╧═╧═╧═╝} \\
\texttt{```} \\
\texttt{Bħalissa qed tilagħbu r{-}round 9 minn massimu ta' 12.} \\
\texttt{Intom flimkien akkumulajtu 5 minn 8 penali.} \\
\texttt{Il{-}plejer l{-}ieħor bagħat dan il{-}messaġġ:} \\
\\ 
\\ 
\texttt{"L huwa f' (3,1) b'suċċess! L{-}oġġetti kollha tiegħi huma mqiegħda: C f' (3,4), P f' (3,5), u L f' (3,1). Inti għandek L f' (3,1) wkoll? Jekk iva, lestejt?"} \\
\\ 
\\ 
\texttt{X'inhu l{-}kmand li jmiss tiegħek?} \\
            }
        }
    }
    & \\ \\

    \theutterance \stepcounter{utterance}  
    & \multicolumn{4}{p{0.6\linewidth}}{
        \cellcolor[rgb]{0.9,0.9,0.9}{
            \makecell[{{p{\linewidth}}}]{
                \texttt{\tiny{[P1$\rangle$GM]}}
                \texttt{GĦID: lestejt!} \\
            }
        }
    }
    & & \\ \\

    \theutterance \stepcounter{utterance}  
    & & & \multicolumn{2}{p{0.3\linewidth}}{
        \cellcolor[rgb]{0.9,0.9,0.9}{
            \makecell[{{p{\linewidth}}}]{
                \texttt{\tiny{[GM$|$GM]}}
                \texttt{GĦID: lestejt!} \\
            }
        }
    }
    & & \\ \\

    \theutterance \stepcounter{utterance}  
    & & & \multicolumn{2}{p{0.3\linewidth}}{
        \cellcolor[rgb]{0.9,0.9,0.9}{
            \makecell[{{p{\linewidth}}}]{
                \texttt{\tiny{[GM$|$GM]}}
                \texttt{true} \\
            }
        }
    }
    & & \\ \\

    \theutterance \stepcounter{utterance}  
    & & & \multicolumn{2}{p{0.3\linewidth}}{
        \cellcolor[rgb]{0.9,0.9,0.9}{
            \makecell[{{p{\linewidth}}}]{
                \texttt{\tiny{[GM$|$GM]}}
                \texttt{Initial states:} \\
\texttt{```} \\
\texttt{Player 1:} \\
\texttt{ 1234567} \\
\texttt{╔═══╤═╤═╗} \\
\texttt{║◌◌◌│◌│◌║ 1} \\
\texttt{║◌P◌└─┤◌║ 2} \\
\texttt{║◌◌◌C◌├─╢ 3} \\
\texttt{╟─┐◌┌─┤◌║ 4} \\
\texttt{║◌│◌│◌│◌║ 5} \\
\texttt{╟─┼─┼─┼─╢ 6} \\
\texttt{║◌│◌│◌│L║ 7} \\
\texttt{╚═╧═╧═╧═╝} \\
\\ 
\\ 
\texttt{Player 2:} \\
\texttt{ 1234567} \\
\texttt{╔═══╤═╤═╗} \\
\texttt{║◌◌◌│◌│◌║ 1} \\
\texttt{║◌◌L└─┤◌║ 2} \\
\texttt{║◌◌◌◌◌├─╢ 3} \\
\texttt{╟─┐◌┌─┤◌║ 4} \\
\texttt{║C│◌│◌│◌║ 5} \\
\texttt{╟─┼─┼─┼─╢ 6} \\
\texttt{║P│◌│◌│◌║ 7} \\
\texttt{╚═╧═╧═╧═╝} \\
\\ 
\texttt{```} \\
            }
        }
    }
    & & \\ \\

    \theutterance \stepcounter{utterance}  
    & & & \multicolumn{2}{p{0.3\linewidth}}{
        \cellcolor[rgb]{0.9,0.9,0.9}{
            \makecell[{{p{\linewidth}}}]{
                \texttt{\tiny{[GM$|$GM]}}
                \texttt{End states:} \\
\texttt{```} \\
\texttt{Player 1:} \\
\texttt{ 1234567} \\
\texttt{╔═══╤═╤═╗} \\
\texttt{║◌◌L│◌│◌║ 1} \\
\texttt{║◌◌◌└─┤◌║ 2} \\
\texttt{║◌◌◌◌◌├─╢ 3} \\
\texttt{╟─┐C┌─┤◌║ 4} \\
\texttt{║◌│P│◌│◌║ 5} \\
\texttt{╟─┼─┼─┼─╢ 6} \\
\texttt{║◌│◌│◌│◌║ 7} \\
\texttt{╚═╧═╧═╧═╝} \\
\\ 
\\ 
\texttt{Player 2:} \\
\texttt{ 1234567} \\
\texttt{╔═══╤═╤═╗} \\
\texttt{║◌◌L│◌│◌║ 1} \\
\texttt{║◌◌◌└─┤◌║ 2} \\
\texttt{║◌◌◌◌◌├─╢ 3} \\
\texttt{╟─┐C┌─┤◌║ 4} \\
\texttt{║◌│P│◌│◌║ 5} \\
\texttt{╟─┼─┼─┼─╢ 6} \\
\texttt{║◌│◌│◌│◌║ 7} \\
\texttt{╚═╧═╧═╧═╝} \\
\\ 
\texttt{```} \\
            }
        }
    }
    & & \\ \\

    \theutterance \stepcounter{utterance}  
    & & & \multicolumn{2}{p{0.3\linewidth}}{
        \cellcolor[rgb]{0.9,0.9,0.9}{
            \makecell[{{p{\linewidth}}}]{
                \texttt{\tiny{[GM$|$GM]}}
                \texttt{* success: True} \\
\texttt{* lose: False} \\
\texttt{* aborted: False} \\
\texttt{{-}{-}{-}{-}{-}{-}{-}} \\
\texttt{* Shifts: 2.00} \\
\texttt{* Max Shifts: 4.00} \\
\texttt{* Min Shifts: 2.00} \\
\texttt{* End Distance Sum: 0.00} \\
\texttt{* Init Distance Sum: 15.11} \\
\texttt{* Expected Distance Sum: 12.57} \\
\texttt{* Penalties: 5.00} \\
\texttt{* Max Penalties: 8.00} \\
\texttt{* Rounds: 9.00} \\
\texttt{* Max Rounds: 12.00} \\
\texttt{* Object Count: 3.00} \\
            }
        }
    }
    & & \\ \\

\end{supertabular}
}

\end{document}
